% Built on the Springer LLNCS template (llncs.cls). Compile with XeLaTeX
% (required for the xeCJK Chinese glyphs).
%
\documentclass[runningheads]{llncs}
\usepackage[T1]{fontenc}
%
\usepackage{amsmath}
\usepackage{amssymb}
\usepackage{graphicx}
\usepackage{booktabs}
\usepackage{array}
\usepackage{xcolor}
\newcommand{\changed}[1]{#1}
\usepackage{tikz}
\usetikzlibrary{arrows.meta,positioning,fit,backgrounds,calc,shapes.geometric}
%
% - Chinese glyph support (XeLaTeX + Noto CJK) -----
\usepackage{xeCJK}
%\usepackage{ctex}

%\setCJKmainfont{Noto Serif CJK SC}
\newcommand{\hz}[1]{{\CJKfamily{}#1}} % inline hanzi helper
% -------------
%

\usepackage{cite}
\usepackage[hidelinks]{hyperref}
% \renewcommand\UrlFont{\ttfamily}
\usepackage{color}
\renewcommand\UrlFont{\color{blue}\rmfamily}
% \urlstyle{rm}

%
% - convenience macros ----------
\newcommand{\cossim}{\sigma}                     % cosine similarity
\newcommand{\jac}{J}                             % Jaccard similarity
\newcommand{\ind}{\mathbb{I}}                    % indicator of a shared trait
\newcommand{\hvec}{\mathbf{h}}                   % hidden state
\newcommand{\avec}{\mathbf{a}}                   % acoustic (pinyin) embedding
\newcommand{\vvec}{\mathbf{v}}                   % semantic (BERT) embedding
\newcommand{\wset}{\mathbf{w}}                   % wubi token set
\newcommand{\pset}{\mathbf{p}}                   % pinyin decomposition
\newcommand{\SCE}{\Sigma_{\mathrm{CE}}}
% -------------
%
\begin{document}
%
\title{Learning Mandarin Phonetic Representations in Neural Networks Using Acoustic Alignment}
\titlerunning{Learning Mandarin Phonetic Representations Using Acoustic Alignment}

\author{Ilya Grigorev}
\authorrunning{I.\,Grigorev}

\institute{Innopolis University, Innopolis, Russia\\
\email{il.grigorev@innopolis.university}}

\maketitle

\section{Motivation}
\label{sec:intro}

Phonetics examines the production, acoustic properties, and perception of speech
sounds \cite{ref_phonetic_sciences}. Different speech sounds may exhibit varying
degrees of phonetic similarity, with some pronunciations being closer to each
other than others. Experimental studies demonstrate that human listeners
distinguish such differences when comparing spoken forms and individual phonemes
\cite{ref_word_similarity,ref_phoneme_similarity}. Graded phonetic similarity
therefore describes an aspect of phonetic structure that extends beyond the
identification of individual sounds. Modeling phonetic information requires
investigating not only whether a system correctly identifies pronunciations but
also whether its learned representations reflect the relationships between
different pronunciations.

We consider the learning of phonetic representations from written Chinese
character sequences, drawing inspiration from the relationship between
orthographic and phonological processing in human reading. Phonological
information is activated during Chinese character recognition, even when readers
perform tasks involving semantic judgments
\cite{ref_early_phonology,ref_chinese_reading}. Mandarin provides a convenient
setting for character-level pronunciation modeling, as each written character
typically corresponds to a single spoken syllable in a given context. Syllables
can be further decomposed into three phonological components: an initial, a
final, and a lexical tone \cite{ref_mandarin_phonology}. Different syllables may
share one or more phonological components, whereas distinct characters may have
identical pronunciations, resulting in substantial \textit{homophony}.
Furthermore, some characters have multiple pronunciations whose selection
depends on the surrounding linguistic context. Learning pronunciation from
written characters therefore requires a model to extract phonetic information
from inputs that also carry orthographic, semantic, and contextual information.
These linguistic properties may influence the organization of learned
representations independently of phonetic similarity. The central challenge is
consequently to determine whether a neural network can acquire representations
organized according to phonetic relationships despite the presence of other
linguistic information in its inputs.


The mapping from character-based representations to pronunciation can be learned
using different forms of phonetic supervision. To encode target pronunciations,
we adopt Hanyu Pinyin (hereafter Pinyin), a standardized Romanization system for
Mandarin Chinese \cite{ref_pinyin}. Rather than predicting complete romanized
syllables, we treat their initial, final, and tone components as separate
categorical labels. Let $y_{\text{ini}}$, $y_{\text{fin}}$, and
$y_{\text{tone}}$ denote the ground-truth component labels, and let
$p_{\text{ini}}$, $p_{\text{fin}}$, and $p_{\text{tone}}$ denote their
respective predicted probability distributions. The discrete prediction loss is
defined as the sum of categorical cross-entropy terms in
\eqref{eq:discrete-loss-component}:

\begin{equation}
\label{eq:discrete-loss-component}
\begin{aligned}
    \mathcal{L}_{\text{pinyin}}
    ={}& -\log p_{\text{ini}}(y_{\text{ini}}) \\
       & -\log p_{\text{fin}}(y_{\text{fin}}) \\
       & -\log p_{\text{tone}}(y_{\text{tone}}).
\end{aligned}
\end{equation}

Supervised classification aims to predict these labels without explicitly
specifying the geometry of the learned representation space. Although shared
phonological components may induce relationships between representations,
accurate classification does not necessarily imply the preservation of graded
similarity between pronunciations. An alternative approach introduces
\textit{continuous acoustic representations} as training targets to guide the
organization of the learned space. Let $\mathbf{h}, \mathbf{a} \in
\mathbb{R}^{d}$ denote the learned hidden representation and its corresponding
acoustic target, respectively. The acoustic alignment loss is defined as the
mean squared error in \eqref{eq:acoustic-loss-component}:

\begin{equation}
\label{eq:acoustic-loss-component}
    \mathcal{L}_{\text{acoustic}}
    = \frac{1}{d}\|\mathbf{h}-\mathbf{a}\|_2^2.
\end{equation}

One approach to obtaining continuous acoustic representations is to extract
embeddings from pretrained speech models. For example, wav2vec 2.0 learns latent
representations from speech waveforms \cite{ref_wav2vec2}. The resulting
acoustic representations encode speech information as continuous vectors rather
than categorical identities, allowing relationships between pronunciations to be
expressed through distances in the acoustic space. However, speech-derived
representations may encode additional properties that are not directly related
to the phonetic distinctions under investigation. Consequently, acoustic
alignment provides an additional source of supervision but does not establish
that the learned geometry corresponds to human-perceived phonetic similarity.

For independent evaluation, we employ DIMSIM, a Mandarin phonetic similarity
framework \cite{ref_dimsim}. DIMSIM represents Pinyin syllables through their
initial, final, and tone components and calculates phonetic distance by
combining distances between the corresponding components. Distances between
initials and finals are obtained from learned multidimensional representations
informed by linguistic structure and native-speaker similarity annotations,
while tonal differences are incorporated through a separate distance component.
The resulting distance measure captures graded relationships between Mandarin
pronunciations and provides a reference independent of the acoustic targets used
during training. By comparing DIMSIM distances with pairwise hidden-state
similarities, we assess whether the learned representations preserve relative
phonetic relationships beyond categorical identity. Although DIMSIM does not
constitute an absolute measure of human-perceived phonetic similarity, the
framework provides a perceptually informed reference for evaluating phonetic
relationships in the learned representation space.

Overall, we compare the representation spaces learned under two principal
training regimes: discrete label supervision \eqref{eq:discrete-loss}, where the
model predicts Pinyin components, and guided supervision
\eqref{eq:combined-loss}, where discrete prediction is supplemented with
continuous speech-derived targets. To investigate the interaction between the
training objectives, we vary the strength of acoustic alignment $\lambda$ from
the label-only regime toward an acoustic-dominated regime, with acoustic-only
training considered as a separate reference condition \eqref{eq:acoustic-loss}. 

\begin{equation}
\label{eq:discrete-loss}
    \mathcal{L} = \mathcal{L}_{\text{pinyin}}.
\end{equation}

\begin{equation}
\label{eq:combined-loss}
    \mathcal{L}
    = \mathcal{L}_{\text{pinyin}}
    + \lambda\mathcal{L}_{\text{acoustic}}.
\end{equation}

\begin{equation}
\label{eq:acoustic-loss}
    \mathcal{L} = \mathcal{L}_{\text{acoustic}}.
\end{equation}

The primary objective is to assess whether graded phonetic similarity emerges
under discrete supervision and how continuous acoustic guidance influences its
preservation in character-based representations.


\section{Expected Results}
\label{sec:results}

The comparison of training regimes allows for several possible outcomes, each
with different implications for the acquisition of phonetic structure. Under
discrete Pinyin supervision, the presence of graded phonetic relationships would
suggest that phonetic similarity can emerge without explicit similarity-based
supervision, whereas their absence would indicate that accurate label prediction
does not necessarily imply the preservation of phonetic similarity. Introducing
acoustic alignment may strengthen the correspondence between hidden-state
similarity and DIMSIM phonetic similarity, although the correspondence may also
remain unchanged, decrease, or vary non-monotonically as the alignment strength
increases. The acoustic-only regime further allows us to examine whether
removing discrete supervision strengthens, preserves, or weakens graded phonetic
relationships compared to the combined and discrete-only regimes. Together, the
possible outcomes provide insight into how discrete pronunciation prediction and
continuous acoustic guidance influence the organization of phonetic information
within character-based representations.


\begin{thebibliography}{10}

\bibitem{ref_phonetic_sciences}
W. J. Hardcastle, J. Laver, and F. E. Gibbon, Eds., \emph{The Handbook of
Phonetic Sciences}, 2nd ed. Chichester, U.K.: Wiley-Blackwell, 2010, doi:
10.1002/9781444317251.

\bibitem{ref_word_similarity}
U. Hahn and T. M. Bailey, ``What makes words sound similar?'' \emph{Cognition},
vol. 97, no. 3, pp. 227--267, 2005, doi: 10.1016/j.cognition.2004.09.006.

\bibitem{ref_phoneme_similarity}
T. M. Bailey and U. Hahn, ``Phoneme similarity and confusability,'' \emph{J.
Memory Lang.}, vol. 52, no. 3, pp. 339--362, 2005, doi:
10.1016/j.jml.2004.12.003.

\bibitem{ref_early_phonology}
C. A. Perfetti and S. Zhang, ``Very early phonological activation in Chinese
reading,'' \emph{J. Exp. Psychol. Learn. Mem. Cogn.}, vol. 21, no. 1, pp.
24--33, 1995, doi: 10.1037/0278-7393.21.1.24.

\bibitem{ref_chinese_reading}
X. Zhou and W. Marslen-Wilson, ``Phonology, orthography, and semantic activation
in reading Chinese,'' \emph{J. Memory Lang.}, vol. 41, no. 4, pp. 579--606,
1999, doi: 10.1006/jmla.1999.2663.

\bibitem{ref_mandarin_phonology}
K. D. Neergaard and C.-R. Huang, ``Mandarin Chinese syllable structure and
phonological similarity: Perception and production studies,'' in \emph{The
Cambridge Handbook of Chinese Linguistics}, C.-R. Huang, Y.-H. Lin, I.-H. Chen,
and Y.-Y. Hsu, Eds. Cambridge, U.K.: Cambridge Univ. Press, 2022, pp. 245--274,
doi: 10.1017/9781108329019.015.

\bibitem{ref_pinyin}
Y.-H. Lin, \emph{The Sounds of Chinese}. Cambridge, U.K.: Cambridge Univ. Press,
2007, ch. 6, pp. 123--135.

\bibitem{ref_wav2vec2}
A. Baevski, H. Zhou, A. Mohamed, and M. Auli, ``wav2vec 2.0: A framework for
self-supervised learning of speech representations,'' in \emph{Advances in
Neural Information Processing Systems}, vol. 33, 2020, pp. 12449--12460.

\bibitem{ref_dimsim}
M. Li, M. Danilevsky, S. Noeman, and Y. Li, ``DIMSIM: An accurate Chinese
phonetic similarity algorithm based on learned high dimensional encoding,'' in
\emph{Proc. 22nd Conf. Computational Natural Language Learning (CoNLL)},
Brussels, Belgium, Oct. 2018, pp. 444--453, doi: 10.18653/v1/K18-1043.

\end{thebibliography}


\end{document}